% Options for packages loaded elsewhere
\PassOptionsToPackage{unicode}{hyperref}
\PassOptionsToPackage{hyphens}{url}
\documentclass[
]{article}
\usepackage{xcolor}
\usepackage[margin=1in]{geometry}
\usepackage{amsmath,amssymb}
\setcounter{secnumdepth}{-\maxdimen} % remove section numbering
\usepackage{iftex}
\ifPDFTeX
  \usepackage[T1]{fontenc}
  \usepackage[utf8]{inputenc}
  \usepackage{textcomp} % provide euro and other symbols
\else % if luatex or xetex
  \usepackage{unicode-math} % this also loads fontspec
  \defaultfontfeatures{Scale=MatchLowercase}
  \defaultfontfeatures[\rmfamily]{Ligatures=TeX,Scale=1}
\fi
\usepackage{lmodern}
\ifPDFTeX\else
  % xetex/luatex font selection
\fi
% Use upquote if available, for straight quotes in verbatim environments
\IfFileExists{upquote.sty}{\usepackage{upquote}}{}
\IfFileExists{microtype.sty}{% use microtype if available
  \usepackage[]{microtype}
  \UseMicrotypeSet[protrusion]{basicmath} % disable protrusion for tt fonts
}{}
\makeatletter
\@ifundefined{KOMAClassName}{% if non-KOMA class
  \IfFileExists{parskip.sty}{%
    \usepackage{parskip}
  }{% else
    \setlength{\parindent}{0pt}
    \setlength{\parskip}{6pt plus 2pt minus 1pt}}
}{% if KOMA class
  \KOMAoptions{parskip=half}}
\makeatother
\usepackage{color}
\usepackage{fancyvrb}
\newcommand{\VerbBar}{|}
\newcommand{\VERB}{\Verb[commandchars=\\\{\}]}
\DefineVerbatimEnvironment{Highlighting}{Verbatim}{commandchars=\\\{\}}
% Add ',fontsize=\small' for more characters per line
\newenvironment{Shaded}{}{}
\newcommand{\AlertTok}[1]{\textcolor[rgb]{1.00,0.00,0.00}{\textbf{#1}}}
\newcommand{\AnnotationTok}[1]{\textcolor[rgb]{0.38,0.63,0.69}{\textbf{\textit{#1}}}}
\newcommand{\AttributeTok}[1]{\textcolor[rgb]{0.49,0.56,0.16}{#1}}
\newcommand{\BaseNTok}[1]{\textcolor[rgb]{0.25,0.63,0.44}{#1}}
\newcommand{\BuiltInTok}[1]{\textcolor[rgb]{0.00,0.50,0.00}{#1}}
\newcommand{\CharTok}[1]{\textcolor[rgb]{0.25,0.44,0.63}{#1}}
\newcommand{\CommentTok}[1]{\textcolor[rgb]{0.38,0.63,0.69}{\textit{#1}}}
\newcommand{\CommentVarTok}[1]{\textcolor[rgb]{0.38,0.63,0.69}{\textbf{\textit{#1}}}}
\newcommand{\ConstantTok}[1]{\textcolor[rgb]{0.53,0.00,0.00}{#1}}
\newcommand{\ControlFlowTok}[1]{\textcolor[rgb]{0.00,0.44,0.13}{\textbf{#1}}}
\newcommand{\DataTypeTok}[1]{\textcolor[rgb]{0.56,0.13,0.00}{#1}}
\newcommand{\DecValTok}[1]{\textcolor[rgb]{0.25,0.63,0.44}{#1}}
\newcommand{\DocumentationTok}[1]{\textcolor[rgb]{0.73,0.13,0.13}{\textit{#1}}}
\newcommand{\ErrorTok}[1]{\textcolor[rgb]{1.00,0.00,0.00}{\textbf{#1}}}
\newcommand{\ExtensionTok}[1]{#1}
\newcommand{\FloatTok}[1]{\textcolor[rgb]{0.25,0.63,0.44}{#1}}
\newcommand{\FunctionTok}[1]{\textcolor[rgb]{0.02,0.16,0.49}{#1}}
\newcommand{\ImportTok}[1]{\textcolor[rgb]{0.00,0.50,0.00}{\textbf{#1}}}
\newcommand{\InformationTok}[1]{\textcolor[rgb]{0.38,0.63,0.69}{\textbf{\textit{#1}}}}
\newcommand{\KeywordTok}[1]{\textcolor[rgb]{0.00,0.44,0.13}{\textbf{#1}}}
\newcommand{\NormalTok}[1]{#1}
\newcommand{\OperatorTok}[1]{\textcolor[rgb]{0.40,0.40,0.40}{#1}}
\newcommand{\OtherTok}[1]{\textcolor[rgb]{0.00,0.44,0.13}{#1}}
\newcommand{\PreprocessorTok}[1]{\textcolor[rgb]{0.74,0.48,0.00}{#1}}
\newcommand{\RegionMarkerTok}[1]{#1}
\newcommand{\SpecialCharTok}[1]{\textcolor[rgb]{0.25,0.44,0.63}{#1}}
\newcommand{\SpecialStringTok}[1]{\textcolor[rgb]{0.73,0.40,0.53}{#1}}
\newcommand{\StringTok}[1]{\textcolor[rgb]{0.25,0.44,0.63}{#1}}
\newcommand{\VariableTok}[1]{\textcolor[rgb]{0.10,0.09,0.49}{#1}}
\newcommand{\VerbatimStringTok}[1]{\textcolor[rgb]{0.25,0.44,0.63}{#1}}
\newcommand{\WarningTok}[1]{\textcolor[rgb]{0.38,0.63,0.69}{\textbf{\textit{#1}}}}
\usepackage{longtable,booktabs,array}
\usepackage{calc} % for calculating minipage widths
% Correct order of tables after \paragraph or \subparagraph
\usepackage{etoolbox}
\makeatletter
\patchcmd\longtable{\par}{\if@noskipsec\mbox{}\fi\par}{}{}
\makeatother
% Allow footnotes in longtable head/foot
\IfFileExists{footnotehyper.sty}{\usepackage{footnotehyper}}{\usepackage{footnote}}
\makesavenoteenv{longtable}
% definitions for citeproc citations
\NewDocumentCommand\citeproctext{}{}
\NewDocumentCommand\citeproc{mm}{%
  \begingroup\def\citeproctext{#2}\cite{#1}\endgroup}
\makeatletter
 % allow citations to break across lines
 \let\@cite@ofmt\@firstofone
 % avoid brackets around text for \cite:
 \def\@biblabel#1{}
 \def\@cite#1#2{{#1\if@tempswa , #2\fi}}
\makeatother
\newlength{\cslhangindent}
\setlength{\cslhangindent}{1.5em}
\newlength{\csllabelwidth}
\setlength{\csllabelwidth}{3em}
\newenvironment{CSLReferences}[2] % #1 hanging-indent, #2 entry-spacing
 {\begin{list}{}{%
  \setlength{\itemindent}{0pt}
  \setlength{\leftmargin}{0pt}
  \setlength{\parsep}{0pt}
  % turn on hanging indent if param 1 is 1
  \ifodd #1
   \setlength{\leftmargin}{\cslhangindent}
   \setlength{\itemindent}{-1\cslhangindent}
  \fi
  % set entry spacing
  \setlength{\itemsep}{#2\baselineskip}}}
 {\end{list}}
\usepackage{calc}
\newcommand{\CSLBlock}[1]{\hfill\break\parbox[t]{\linewidth}{\strut\ignorespaces#1\strut}}
\newcommand{\CSLLeftMargin}[1]{\parbox[t]{\csllabelwidth}{\strut#1\strut}}
\newcommand{\CSLRightInline}[1]{\parbox[t]{\linewidth - \csllabelwidth}{\strut#1\strut}}
\newcommand{\CSLIndent}[1]{\hspace{\cslhangindent}#1}
\setlength{\emergencystretch}{3em} % prevent overfull lines
\providecommand{\tightlist}{%
  \setlength{\itemsep}{0pt}\setlength{\parskip}{0pt}}
\usepackage{bookmark}
\IfFileExists{xurl.sty}{\usepackage{xurl}}{} % add URL line breaks if available
\urlstyle{same}
\hypersetup{
  pdftitle={Specification Surfaces as Source-of-Truth Artifacts in AI-Assisted Software Engineering},
  pdfauthor={Petros Lambropoulos (Rooted Layers), petroslamb.dev@gmail.com},
  hidelinks,
  pdfcreator={LaTeX via pandoc}}

\title{Specification Surfaces as Source-of-Truth Artifacts in
AI-Assisted Software Engineering}
\author{Petros Lambropoulos (Rooted Layers), petroslamb.dev@gmail.com}
\date{May 2026}

\begin{document}
\maketitle
\begin{abstract}
AI-assisted software review usually begins with the generated diff, even
when the decision a model followed lives somewhere else: a prompt, chat
transcript, stale note, workflow file, test, harness, approval record,
or private design choice. That mismatch creates a source-of-truth
problem. Reviewers can inspect the downstream artifact while the
artifact that authorized, constrained, or evaluated the work may remain
outside the governed system. This paper defines a
\texttt{specification\ surface} as a maintained artifact with
operational standing in AI-assisted software work because people or
tools use it to constrain, validate, trace, authorize, or repair
generated behavior. The definition includes specifications, schemas,
tests, workflow contracts, harnesses, trace records, claim ledgers, and
approval gates when they are maintained and coupled to control or
evaluation. It excludes passive documentation and one-off prompts. The
contribution is an eight-role taxonomy, a lifecycle model of
\texttt{specify\ -\textgreater{}\ generate\ -\textgreater{}\ validate\ -\textgreater{}\ trace\ -\textgreater{}\ repair},
a systematization of recent AI-assisted software-engineering literature,
and an illustrative artifact analysis of a file-backed workflow packet.
The case is not evidence that one workflow improves outcomes. It shows
how the framework separates source authority from evidence, traces, and
approval records, giving code review a way to ask which upstream
artifacts governed the code it is asked to judge.
\end{abstract}

\section{Keywords}\label{keywords}

AI-assisted software engineering; specification; code review; agent
workflows; traceability; software process; governance.

\section{1. Introduction}\label{introduction}

Review starts at the diff because the diff is visible. In AI-assisted
software work, that habit is increasingly incomplete. A model may have
generated the patch from an instruction file, a task contract, a test
scaffold, a harness, a workflow policy, or a conversation that never
becomes part of the repository. The implementation can look ordinary
while the decision surface that governed it is missing, stale, or
unreviewed.

Generated code still deserves ordinary engineering scrutiny. It can be
wrong, insecure, slow, brittle, or inconsistent with local design. The
problem is that the diff may not reveal the artifact that made the
change reasonable or unreasonable in the first place. A reviewer can ask
whether the code compiles and whether the tests pass, yet still be
unable to inspect the maintained statement of intended behavior, the
constraints the agent was expected to preserve, the validation surface
that checked those constraints, or the repair path that should update
the upstream artifact after failure.

This is a source-of-truth problem rather than a claim about syntax. The
question is not whether natural language is code. The question is which
artifacts have enough operational standing to govern AI-assisted
software work. A design note, prompt, README, workflow file, or test
suite can all influence a model. Only some of them deserve to be treated
as source-of-truth surfaces.

This paper uses \texttt{specification\ surface} for that narrower class:
maintained artifacts that become part of the source-of-truth boundary
because people or tools use them to constrain, validate, trace,
authorize, or repair generated behavior. A surface may be written in
prose, YAML, a schema language, tests, a workflow DSL, a trace format,
or a review checklist. File format is secondary. The defining property
is operational coupling.

Operational coupling also narrows the role of natural language. A
natural-language instruction does not become a governing specification
merely because a harness can interpret it. It becomes part of a
governing surface only when the surrounding harness binds it to
executable lifecycle semantics: artifact contracts, state, adapters,
validation, and recovery behavior. Otherwise it remains an instruction
or interface layer.

The paper makes four contributions:

\begin{enumerate}
\def\labelenumi{\arabic{enumi}.}
\tightlist
\item
  a definition of \texttt{specification\ surface} for AI-assisted
  software engineering;
\item
  an eight-role taxonomy covering source, runtime, validation, trace,
  evidence, context, repair, and authority surfaces;
\item
  a lifecycle model,
  \texttt{specify\ -\textgreater{}\ generate\ -\textgreater{}\ validate\ -\textgreater{}\ trace\ -\textgreater{}\ repair},
  that explains why failures should update upstream surfaces rather than
  only downstream code;
\item
  a literature mapping and an illustrative artifact analysis showing how
  the framework classifies existing systems and one file-backed workflow
  packet.
\end{enumerate}

The paper does not argue for maximal formalism. It argues for
identifying which artifacts actually govern the work, how they are
checked, and when they should be repaired.

\section{2. Background And Motivation}\label{background-and-motivation}

\subsection{2.1 Specification authority in software
work}\label{specification-authority-in-software-work}

Spec-driven development treats specifications as active development
artifacts rather than background requirements prose. Piskala describes a
spectrum from spec-first work, to spec-anchored work, to spec-as-source
generation, with stronger specification discipline justified when it
removes consequential ambiguity and can be maintained
(\citeproc{ref-piskala2026specdriven}{Piskala 2026}). That spectrum
gives the present paper a software-engineering baseline. Specification
surfaces are not limited to agent harnesses or prompt files; they
continue a longer concern that software becomes hard to govern when
intended behavior is not itself maintained.

The framework here is narrower than a general recommendation to
formalize everything. A surface earns authority when it participates in
control, validation, traceability, or repair. A highly formal artifact
that nobody maintains can be weaker than a short acceptance contract
that drives generation and tests.

\subsection{2.2 Specification gaps in agentic
generation}\label{specification-gaps-in-agentic-generation}

The issue becomes sharper when agents coordinate under partial
knowledge. Sartori studies multi-agent class generation while reducing
specification detail and introducing opposing representation biases
(\citeproc{ref-sartori2026specgap}{Sartori 2026}). The useful lesson is
not that every task needs maximal detail. It is that underspecified
shared intent can become an integration failure, and that restoring
specification can be more useful than post-hoc conflict information in
the studied setting.

That result anchors the review problem. When a generated implementation
diverges, the missing object may sit upstream of the diff: the shared
decision that agents needed but did not receive. A review process that
only inspects output can miss the ungoverned decision surface that
caused the output to look plausible and still be wrong.

\subsection{2.3 Professional control around
agents}\label{professional-control-around-agents}

Huang et al.~report that experienced developers use coding agents
through planning, supervision, validation, version control, and domain
expertise (\citeproc{ref-huang2025professional}{Huang et al. 2025}).
That practice matters because it locates AI assistance inside
professional control rather than outside it. The human task is not
merely to accept or reject a model output. It is to preserve the
artifacts that make design intent inspectable and failure repairable.

The opposite pattern is informal delegation: intent remains in prompt
material or conversation, while generated code becomes the only artifact
reviewed with seriousness. The technical point is that some artifacts
need to carry authority, some only provide context, and some should be
repaired when failure exposes missing intent.

\subsection{2.4 Method and scope}\label{method-and-scope}

The framework was built from three inputs. First, prior
software-engineering work on spec-driven development, agent
specification gaps, and professional control established the core
problem: intended behavior becomes difficult to govern when
specification authority is absent, stale, or unreviewed. Second, recent
agent and validation systems were coded by artifact role and lifecycle
stage to test whether the surface vocabulary separates source, runtime,
validation, trace, evidence, repair, context, and authority functions.
Third, one local file-backed workflow packet was analyzed as an
illustrative case.

This is a framework and systematization article, not a systematic
literature review. The literature corpus is targeted around the
source-of-truth problem in AI-assisted software work. The case study is
intentionally modest: it tests whether the taxonomy can classify a real
artifact packet and expose authority/evidence/trace distinctions. It
does not measure effectiveness.

\section{3. Definition: Specification
Surface}\label{definition-specification-surface}

A specification surface is a maintained artifact that becomes part of
the source-of-truth boundary for AI-assisted software work because
people or tools use it to constrain, validate, trace, authorize, or
repair generated behavior.

An artifact qualifies when at least one of the following is true:

\begin{itemize}
\tightlist
\item
  it constrains allowed behavior before execution;
\item
  tooling parses, types, lints, interprets, compiles, or validates it;
\item
  it drives generation, orchestration, decomposition, or tool access;
\item
  it defines tests, checks, gates, policies, or completion criteria;
\item
  it links outputs to governing intent or supporting evidence;
\item
  it records execution state that can be inspected, replayed, or
  compared;
\item
  failure turns the artifact itself into a repair target.
\end{itemize}

This definition excludes passive documentation. A README does not become
a specification surface merely because an agent can read it. A prompt
does not become source of truth because it came before generated code. A
context file becomes a candidate surface only when it is maintained and
connected to evaluation, enforcement, runtime interpretation,
traceability, or review.

The boundary also protects four claims that often get blurred in
AI-assisted software discussions. Natural language can participate in a
governed surface, but natural language is not automatically code. In the
Natural-Language Agent Harnesses setting, the natural-language layer is
not sufficient by itself; IHR-style runtime mechanisms, interfaces, and
execution semantics are what make the harness operational. Documentation
can inform work, but documentation is not automatically authority.
Generated code remains consequential because it is the artifact users
and systems execute. More formal specification can help when it removes
consequential ambiguity and can be checked, but formality is not a
monotonic improvement in model behavior.

\section{4. Taxonomy Of Specification
Surfaces}\label{taxonomy-of-specification-surfaces}

Table 1 treats surface types as roles rather than mutually exclusive
file formats. A single artifact can occupy several roles when it is both
human-reviewable and machine-interpreted.

\textbf{Table 1. Specification-surface taxonomy}

\begin{longtable}[]{@{}
  >{\raggedright\arraybackslash}p{(\linewidth - 8\tabcolsep) * \real{0.2000}}
  >{\raggedright\arraybackslash}p{(\linewidth - 8\tabcolsep) * \real{0.2000}}
  >{\raggedright\arraybackslash}p{(\linewidth - 8\tabcolsep) * \real{0.2000}}
  >{\raggedright\arraybackslash}p{(\linewidth - 8\tabcolsep) * \real{0.2000}}
  >{\raggedright\arraybackslash}p{(\linewidth - 8\tabcolsep) * \real{0.2000}}@{}}
\toprule\noalign{}
\begin{minipage}[b]{\linewidth}\raggedright
Surface type
\end{minipage} & \begin{minipage}[b]{\linewidth}\raggedright
Working definition
\end{minipage} & \begin{minipage}[b]{\linewidth}\raggedright
Operational force
\end{minipage} & \begin{minipage}[b]{\linewidth}\raggedright
Examples
\end{minipage} & \begin{minipage}[b]{\linewidth}\raggedright
Main risks
\end{minipage} \\
\midrule\noalign{}
\endhead
\bottomrule\noalign{}
\endlastfoot
Source surfaces & Artifacts that state desired behavior, constraints,
interfaces, or acceptance criteria before implementation. & Guide
generation, constrain implementation, establish review target. & Product
specs, API contracts, schemas, BDD scenarios, task contracts. & Drift,
ambiguity, over-specification, false confidence when the spec is
wrong. \\
Runtime surfaces & Artifacts interpreted by a runtime to orchestrate
agent or workflow behavior. & Execute control flow, roles, state
transitions, tool access, and decomposition. & Agent harnesses, workflow
YAML, agent configs, DSLs, skill definitions. & Hidden runtime
semantics, weak replayability, model-dependent interpretation. \\
Validation surfaces & Artifacts that check whether behavior conforms to
stated intent. & Turn expectations into tests, validators, linters, or
contract checks. & Test suites, generated checks, static validators,
mutation tests. & Test gaming, incomplete coverage, validator drift. \\
Trace surfaces & Artifacts that retain execution evidence for audit,
replay, debugging, or comparison. & Preserve what ran, which inputs were
admitted, and which constraints applied. & Run logs, agent traces,
checkpoint logs, replay records. & Treating logs as authority rather
than evidence. \\
Evidence surfaces & Artifacts that link claims, outputs, or decisions to
support and contradictions. & Make claims auditable and contradictions
visible. & Claim ledgers, source maps, evidence tables, provenance
notes. & Paper-thin support, citation laundering, stale evidence. \\
Context surfaces & Project guidance that becomes operational only when
maintained and evaluated. & Shape generation or review when coupled to
checks or process. & Repository instructions and coding guidelines. &
Passive context mistaken for enforcement. \\
Repair surfaces & Artifacts that convert failure diagnosis into upstream
revision. & Feed validation or trace failures back into specs, policies,
tests, or prompts. & Repair tickets, revised specs, updated tests,
policy patches. & Fixing only generated output while leaving controlling
ambiguity intact. \\
Authority and gate surfaces & Artifacts that decide what may run, move,
publish, or be promoted. & Govern lifecycle movement, permissions,
promotion, or external mutation. & Gate contracts, approval records,
release checklists, policy manifests. & Bypassed approvals; unclear
authority order. \\
\end{longtable}

The taxonomy is intentionally plural. Source authority may live in a
schema, a workflow contract, a test-generated prompt, or a human
approval record. No single surface should replace all others. Review
needs to know which surfaces govern which decisions.

\section{5. Lifecycle Model}\label{lifecycle-model}

Specification surfaces matter because they participate in a lifecycle
rather than remaining static documents.

\begin{Shaded}
\begin{Highlighting}[]
\NormalTok{specify {-}\textgreater{} generate {-}\textgreater{} validate {-}\textgreater{} trace {-}\textgreater{} repair}
\NormalTok{   \^{}                                      |}
\NormalTok{   |\_\_\_\_\_\_\_\_\_\_\_\_\_\_\_\_\_\_\_\_\_\_\_\_\_\_\_\_\_\_\_\_\_\_\_\_\_\_|}
\end{Highlighting}
\end{Shaded}

In the first stage, a surface states intent, constraints, or authority.
In the generation stage, an agent, tool, or workflow produces code or
another downstream artifact. Validation checks conformance. Trace
records what happened and which inputs or surfaces were admitted. Repair
updates the governing surface when failures expose ambiguity, drift,
stale assumptions, missing constraints, or weak evidence.

This loop separates governed specification surfaces from ordinary
documents. A surface should be reviewable before execution, useful
during execution, and repairable after failure.

Recent systems occupy different parts of this loop. Bootstrapping Coding
Agents treats the specification as the stable artifact of record in a
narrow agent-bootstrap setting
(\citeproc{ref-monperrus2026bootstrapping}{Monperrus 2026}).
Natural-Language Agent Harnesses and AgentSPEX make harness or workflow
logic executable through runtimes with contracts, state, checkpointing,
tracing, replay, or verification (\citeproc{ref-pan2026nlah}{Pan et al.
2026}; \citeproc{ref-wang2026agentspex}{Wang et al. 2026}). For NLAH
specifically, the relevant mechanism is IHR plus NLAH: the
natural-language harness description becomes operational because an
underlying harness or framework supplies interfaces and execution
semantics, not because natural language alone becomes code. ContextCov
converts passive agent instruction files into executable constraints,
illustrating the transition from context surface to validation surface
(\citeproc{ref-sharma2026contextcov}{Sharma 2026}). Test-Driven AI Agent
Definition compiles behavioral specifications into tests and prompt
artifacts (\citeproc{ref-rehan2026tdad}{Rehan 2026}). FASTRIC treats
prompt-protocol formality as a design parameter for verifiable
interactions (\citeproc{ref-jin2025fastric}{Jin 2025}). Trace and repair
work such as AgentPex, AgentFixer, and claim-level auditability shows
why outcome-only review remains weak for long-running agent behavior
(\citeproc{ref-barke2026willful}{Barke et al. 2026};
\citeproc{ref-mulian2026agentfixer}{Mulian et al. 2026};
\citeproc{ref-rasheed2026fluent}{Rasheed et al. 2026}).

The lifecycle model also explains why generated code remains important.
Generated code is the executable artifact that users and systems
encounter. But when it fails, the repair target is often not only the
code. It is also the governing surface that allowed the failure to
happen.

\section{6. Literature Mapping}\label{literature-mapping}

Table 2 maps the core literature by surface role and lifecycle stage.
The table is a systematization of artifact roles in recent AI-assisted
software-engineering work, not a broad survey of the area.

\textbf{Table 2. Literature mapped to specification-surface roles}

\begin{longtable}[]{@{}
  >{\raggedright\arraybackslash}p{(\linewidth - 8\tabcolsep) * \real{0.2000}}
  >{\raggedright\arraybackslash}p{(\linewidth - 8\tabcolsep) * \real{0.2000}}
  >{\raggedright\arraybackslash}p{(\linewidth - 8\tabcolsep) * \real{0.2000}}
  >{\raggedright\arraybackslash}p{(\linewidth - 8\tabcolsep) * \real{0.2000}}
  >{\raggedright\arraybackslash}p{(\linewidth - 8\tabcolsep) * \real{0.2000}}@{}}
\toprule\noalign{}
\begin{minipage}[b]{\linewidth}\raggedright
Source
\end{minipage} & \begin{minipage}[b]{\linewidth}\raggedright
Surface roles
\end{minipage} & \begin{minipage}[b]{\linewidth}\raggedright
Lifecycle stages
\end{minipage} & \begin{minipage}[b]{\linewidth}\raggedright
Use in this paper
\end{minipage} & \begin{minipage}[b]{\linewidth}\raggedright
Boundary
\end{minipage} \\
\midrule\noalign{}
\endhead
\bottomrule\noalign{}
\endlastfoot
Spec-driven development (\citeproc{ref-piskala2026specdriven}{Piskala
2026}) & source, validation, repair & specify, generate, validate,
repair & SE baseline for spec-first, spec-anchored, and spec-as-source
modes. & Do not claim maximum specification is always best. \\
Specification gap (\citeproc{ref-sartori2026specgap}{Sartori 2026}) &
source, validation, repair & specify, generate, validate, repair & Hinge
evidence that underspecification can cause coordination failure. &
Specific tasks and models; do not overgeneralize. \\
Professional control (\citeproc{ref-huang2025professional}{Huang et al.
2025}) & source, validation, repair, authority & specify, validate,
repair & Practice baseline for expert supervision and validation. & Does
not prove universal use of specification surfaces. \\
Bootstrap specs (\citeproc{ref-monperrus2026bootstrapping}{Monperrus
2026}) & source, validation & specify, generate, validate & Narrow
example of specification as record. & Does not make all implementation
safely regenerable. \\
NL agent harnesses (\citeproc{ref-pan2026nlah}{Pan et al. 2026}) &
runtime, source, trace & specify, generate, trace & Natural-language
harness descriptions made operational through IHR-style runtime
semantics, artifact contracts, state, adapters, validation, and recovery
behavior. & Do not treat every harness-interpreted instruction as a
governing specification; without executable lifecycle binding it remains
instruction/interface layer. \\
AgentSPEX (\citeproc{ref-wang2026agentspex}{Wang et al. 2026}) &
runtime, validation, trace, source & specify, generate, validate, trace
& Typed workflow specification and execution language. & Not every team
needs a DSL. \\
ContextCov (\citeproc{ref-sharma2026contextcov}{Sharma 2026}) & context,
validation, repair & specify, validate, repair & Passive instructions
transformed into executable guardrails. & Ordinary context files are not
automatically governed. \\
Test-driven agent definition (\citeproc{ref-rehan2026tdad}{Rehan 2026})
& source, validation, repair & specify, generate, validate, repair &
Behavioral specs compiled into tests and prompt artifacts. & Tests do
not prove all behavior. \\
FASTRIC (\citeproc{ref-jin2025fastric}{Jin 2025}) & source, validation,
trace & specify, validate, trace & Prompt-protocol specification and
procedural conformance. & Formality is not monotonic improvement. \\
AgentPex traces (\citeproc{ref-barke2026willful}{Barke et al. 2026}) &
trace, validation & validate, trace & Procedural trace checking beyond
outcome success. & Traces diagnose; they do not replace source
authority. \\
AgentFixer (\citeproc{ref-mulian2026agentfixer}{Mulian et al. 2026}) &
validation, trace, repair & validate, trace, repair & Failure detection,
root cause, and repair recommendations. & Automated repair does not
replace human review. \\
Claim auditability (\citeproc{ref-rasheed2026fluent}{Rasheed et al.
2026}) & evidence, trace, validation, repair & trace, validate, repair &
Claim-level auditability and contradiction transparency. & This case
does not implement full provenance graphs. \\
\end{longtable}

The mapping supports a bounded synthesis: source-of-truth authority can
sit upstream of generated code, but only when artifacts are connected to
control, validation, traceability, or repair.

\section{7. Illustrative Artifact
Analysis}\label{illustrative-artifact-analysis}

\subsection{7.1 Method}\label{method}

This section presents an illustrative artifact analysis, not a
controlled empirical evaluation. The unit of analysis is a single
specification-surface package in a file-backed editorial and
software-workflow repository, together with the paper-development and
essay-iteration workflow machinery around it.

The artifact corpus includes package status, package README,
claim-boundary ledger, paper blueprint, publication route note, workflow
contracts, gate contracts, policy manifests, retained run logs, external
review synthesis, and human promotion approval records. Each artifact
was coded against the surface roles in Table 1 and the lifecycle stages
in Section 5. The coding protocol asked three questions:

\begin{enumerate}
\def\labelenumi{\arabic{enumi}.}
\tightlist
\item
  Which decision does the artifact govern, if any?
\item
  Which tool, reviewer, or workflow step reads or validates it?
\item
  When a failure occurs, is the artifact evidence, authority, or a
  repair target?
\end{enumerate}

The analysis reports observations only where the artifact role is
visible from maintained files or retained run records. Interpretive
claims are limited to what the taxonomy can classify: source authority,
runtime interpretation, validation, trace, evidence, repair, context, or
gate behavior. The case is useful because it separates roles that often
collapse into a single conversational stream during AI-assisted work. It
is also limited: it is a single local system, analyzed by the author,
and does not prove general effectiveness.

\subsection{7.2 Case summary}\label{case-summary}

Observation 1: the package has a \texttt{status.yaml} manifest that
resolves its kind, revision, live draft, evidence path, and lifecycle
stage. Package state is therefore a maintained surface rather than an
inference from prose.

Observation 2: root authority surfaces define which layer wins when
artifacts disagree. Workflow contracts, policy YAMLs, gate contracts,
pipeline profiles, and package status outrank logs and views. This keeps
trace evidence from silently becoming source authority.

Observation 3: the experimental \texttt{paper\_development} workflow
admits package status, package README, the live draft or candidate,
claim boundary ledger, evidence surface, and paper folder. Its allowed
write path is the package-local \texttt{paper/} folder, keeping paper
work separate from the public essay route.

Observation 4: the completed public-essay route supplies a repair and
promotion example. External reviews identified non-blocking polish
requirements and judged the candidate to clear the north-star bar, but
the external review synthesis did not authorize live draft replacement.
A separate human promotion approval surface authorized copying the
polished candidate into the live draft.

These observations support a limited claim: the taxonomy can classify
source authority, runtime contracts, evidence, trace, and approval as
distinct surfaces in a real workflow. They do not support a claim that
this repository's workflow improves quality, speed, safety, or adoption
compared with another process.

\subsection{7.3 Artifact classification}\label{artifact-classification}

Table 3 uses role labels rather than full repository paths so the
manuscript remains readable in narrow columns. The exact file names are
retained in the local case-study trace and workflow packet.

\textbf{Table 3. Local artifact roles}

\begin{longtable}[]{@{}
  >{\raggedright\arraybackslash}p{(\linewidth - 8\tabcolsep) * \real{0.2000}}
  >{\raggedright\arraybackslash}p{(\linewidth - 8\tabcolsep) * \real{0.2000}}
  >{\raggedright\arraybackslash}p{(\linewidth - 8\tabcolsep) * \real{0.2000}}
  >{\raggedright\arraybackslash}p{(\linewidth - 8\tabcolsep) * \real{0.2000}}
  >{\raggedright\arraybackslash}p{(\linewidth - 8\tabcolsep) * \real{0.2000}}@{}}
\toprule\noalign{}
\begin{minipage}[b]{\linewidth}\raggedright
Artifact
\end{minipage} & \begin{minipage}[b]{\linewidth}\raggedright
Surface type
\end{minipage} & \begin{minipage}[b]{\linewidth}\raggedright
Lifecycle stage
\end{minipage} & \begin{minipage}[b]{\linewidth}\raggedright
Operational force
\end{minipage} & \begin{minipage}[b]{\linewidth}\raggedright
Failure or repair implication
\end{minipage} \\
\midrule\noalign{}
\endhead
\bottomrule\noalign{}
\endlastfoot
Authority map & authority, gate & specify & Defines source-of-truth
order. & If logs act as source authority, repair the authority map or
consuming workflow. \\
Package status manifest & authority, context & specify & Resolves
package kind, revision, live draft, evidence path, and stage state. &
Stage decisions should be repaired here, not inferred from prose
alone. \\
Paper workflow contract & runtime, authority, validation & specify,
validate, trace & Declares context dependencies, allowed paths, outputs,
and checks. & Paper artifacts cannot silently mutate the essay route. \\
Essay workflow contract & runtime, validation, authority & specify,
generate, validate, trace & Declares candidate, review, scorecard,
decision, and gate outputs. & Promotion requires review and approval
surfaces. \\
Gate contract manifest & authority, validation & validate, repair &
Defines movement, quality, promotion, and external-mutation gates. &
Missing approvals should block movement or mutation. \\
Claim boundary ledger & evidence, repair & specify, validate, repair &
Names authorized, caveated, and forbidden claims. & Paper prose that
drifts into forbidden claims repairs the blueprint or draft. \\
External review synthesis & evidence, repair & validate, repair &
Aggregates review findings and required polish. & Review can repair a
candidate without authorizing promotion. \\
Human promotion approval & authority, gate & validate, repair & Records
explicit approval to promote the polished candidate. & Promotion is
authorized through this surface, not through file existence. \\
Reviewable run logs & trace & trace & Retain operational evidence. &
Logs diagnose what happened but are not source authority. \\
\end{longtable}

\subsection{7.4 Case lesson}\label{case-lesson}

The local case shows how a workflow can keep source authority, runtime
interpretation, evidence, trace, and approval distinct. It does not show
that this repository's process is generally better than alternatives.
The value of the case is narrower: the taxonomy can classify a real
file-backed workflow and identify which artifacts govern intent before,
during, and after AI-assisted work.

\section{8. Discussion}\label{discussion}

\subsection{8.1 Code review moves upstream, but code remains
consequential}\label{code-review-moves-upstream-but-code-remains-consequential}

Generated code does not stop mattering. It may simply be the second
artifact reviewers should ask about. When generated output is downstream
of maintained specification surfaces, reviewers need to inspect both the
diff and the surface that governed the diff. A clean diff can still be
wrong if the controlling surface is stale, ambiguous, or unauthorized.

\subsection{8.2 Context files become surfaces only under
control}\label{context-files-become-surfaces-only-under-control}

Agent instruction files and repository context files deserve attention
because they are machine-facing project artifacts. A context file
becomes a governed specification surface only when it is maintained,
evaluated, enforced, interpreted by a runtime, or tied to review.
ContextCov is useful here because it turns passive instructions into
executable guardrails (\citeproc{ref-sharma2026contextcov}{Sharma
2026}). Presence in the context window is weaker than operational
coupling.

\subsection{8.3 Formality is a design
parameter}\label{formality-is-a-design-parameter}

The taxonomy should not be read as a call for maximal formalism. A more
formal surface can help when it removes consequential ambiguity and can
be checked. It can also add friction, encode the wrong behavior, or
over-constrain work when the domain is still exploratory. FASTRIC is
useful here because it treats prompt-protocol formality as a design
variable rather than a moral victory (\citeproc{ref-jin2025fastric}{Jin
2025}).

\subsection{8.4 Failure modes}\label{failure-modes}

Specification surfaces can fail in familiar ways:

\begin{itemize}
\tightlist
\item
  they drift from implementation or practice;
\item
  they become too broad to review;
\item
  they encode stale authority;
\item
  validators check the wrong thing;
\item
  generated tests create false confidence;
\item
  traces are mistaken for source truth;
\item
  approval gates become ceremony rather than decision points.
\end{itemize}

These failures do not invalidate the concept. They explain why the
surface must be maintained, validated, traced, and repaired.

\section{9. Limitations And Threats To
Validity}\label{limitations-and-threats-to-validity}

This is a framework and systematization paper with an illustrative
artifact analysis. It is not a benchmark paper and does not measure
general effectiveness.

The taxonomy is conceptual and should be tested against additional
systems, teams, and workflow architectures. The case study is
single-case, local, and author-coded. It demonstrates how the taxonomy
can classify a file-backed workflow system; it does not prove that the
workflow improves quality, speed, safety, or adoption.

The literature map is preprint-heavy and fast-moving. That is
appropriate for an early framework paper in AI-assisted software
engineering, but it increases the risk that terminology, systems, and
empirical results will change quickly.

The argument does not claim that natural language is code. Some
specification surfaces may be written in ordinary language, but the
defining property is operational coupling, not syntax. Natural-language
harness descriptions count as governing surfaces only when they are
bound to executable lifecycle semantics; otherwise they remain
instruction or interface layers. The paper also does not claim that all
documentation becomes authoritative, that generated code is disposable,
or that richer specification always improves model behavior.

\section{10. Artifact Availability}\label{artifact-availability}

This manuscript analyzes repository-local artifacts from the author's
content workflow. The artifact packet consists of the paper roadmap,
north-star rubric, draftability audit, approval gate, case-study trace,
literature map, manuscript drafts, route checklists, and retained
workflow records. This version does not include a separate public
supplemental bundle. The case analysis is therefore presented as
illustrative and author-audited rather than as an externally reusable
dataset.

The next artifact step is a public or anonymized snapshot that lets
readers inspect the classification method and workflow evidence
directly. The packet is not presented as a benchmark dataset.

\section{11. AI/Tool Use Disclosure}\label{aitool-use-disclosure}

This manuscript was prepared with AI-assisted writing and
repository-analysis tooling under a file-backed workflow. AI assistance
was used to organize local evidence, draft and revise manuscript prose,
and check workflow artifacts against the paper-development rubric. The
author remains responsible for claims, citations, artifact selection,
revision decisions, and public-posting decisions.

\section{12. Conclusion}\label{conclusion}

AI-assisted software work becomes easier to govern when intent can be
read before execution, checked during execution, traced after execution,
and repaired when failures expose missing decisions. Code review remains
necessary, but it is incomplete when the upstream source of authority is
informal. Specification surfaces name the artifacts that make intent
reviewable: source surfaces, runtime surfaces, validation surfaces,
trace surfaces, evidence surfaces, context surfaces, repair surfaces,
and authority gates. The practical question for AI-assisted software
engineering is not whether prompts are code. It is whether the artifacts
that control generated work are maintained enough to be governed.

\section{References}\label{references}

\protect\phantomsection\label{refs}
\begin{CSLReferences}{1}{1}
\bibitem[\citeproctext]{ref-barke2026willful}
Barke, Shraddha, Reshabh K. Sharma, and Benjamin Zorn. 2026.
\emph{Willful Disobedience: Automatically Detecting Failures in Agentic
Traces}. \url{https://arxiv.org/abs/2603.23806}.

\bibitem[\citeproctext]{ref-huang2025professional}
Huang, Ruanqianqian, Avery Reyna, Sorin Lerner, Haijun Xia, and Brian
Hempel. 2025. \emph{Professional Software Developers Don't Vibe, They
Control: AI Agent Use for Coding in 2025}.
\url{https://arxiv.org/abs/2512.14012}.

\bibitem[\citeproctext]{ref-jin2025fastric}
Jin, Wen-Long. 2025. \emph{FASTRIC: Prompt Specification Language for
Verifiable LLM Interactions}. \url{https://arxiv.org/abs/2512.18940}.

\bibitem[\citeproctext]{ref-monperrus2026bootstrapping}
Monperrus, Martin. 2026. \emph{Bootstrapping Coding Agents: The
Specification Is the Program}. \url{https://arxiv.org/abs/2603.17399}.

\bibitem[\citeproctext]{ref-mulian2026agentfixer}
Mulian, Hadar, Liane Galanti, Sergey Zeltyn, Avi Yaeli, Ido Levy, and
Segev Shlomov. 2026. \emph{AgentFixer: From Failure Detection to Fix
Recommendations in LLM Agentic Systems}.
\url{https://arxiv.org/abs/2603.29848}.

\bibitem[\citeproctext]{ref-pan2026nlah}
Pan, Linyue, Lexiao Zou, Shuo Guo, Jingchen Ni, and Hai-Tao Zheng. 2026.
\emph{Natural-Language Agent Harnesses}.
\url{https://arxiv.org/abs/2603.25723}.

\bibitem[\citeproctext]{ref-piskala2026specdriven}
Piskala, Deepak Babu. 2026. \emph{Spec-Driven Development: From Code to
Contract in the Age of AI Coding Assistants}.
\url{https://arxiv.org/abs/2602.00180}.

\bibitem[\citeproctext]{ref-rasheed2026fluent}
Rasheed, Razeen A., Animesh Mukherjee, Somnath Banerjee, and Rima Hazra.
2026. \emph{From Fluent to Verifiable: Claim-Level Auditability for Deep
Research Agents}. \url{https://arxiv.org/abs/2602.13855}.

\bibitem[\citeproctext]{ref-rehan2026tdad}
Rehan, Tzafrir. 2026. \emph{Test-Driven AI Agent Definition: Compiling
Tool-Using Agents from Behavioral Specifications}.
\url{https://arxiv.org/abs/2603.08806}.

\bibitem[\citeproctext]{ref-sartori2026specgap}
Sartori, Camilo Chacón. 2026. \emph{The Specification Gap: Coordination
Failure Under Partial Knowledge in Code Agents}.
\url{https://arxiv.org/abs/2603.24284}.

\bibitem[\citeproctext]{ref-sharma2026contextcov}
Sharma, Reshabh K. 2026. \emph{ContextCov: Deriving and Enforcing
Executable Constraints from Agent Instruction Files}.
\url{https://arxiv.org/abs/2603.00822}.

\bibitem[\citeproctext]{ref-wang2026agentspex}
Wang, Pengcheng, Jerry Huang, Jiarui Yao, et al. 2026. \emph{AgentSPEX:
An Agent Specification and Execution Language}.
\url{https://arxiv.org/abs/2604.13346}.

\end{CSLReferences}

\end{document}
